\section{\texorpdfstring{Helicity Suppression in $\pi\to\ell\nu$}{Helicity Suppression in pi to l nu}}

\paragraph{Negative-helicity antiparticle in the chiral basis}

A negative-helicity antiparticle spinor may be written as
\begin{align}
v_\downarrow
=
\sqrt{E+m}
\begin{pmatrix}
\dfrac{|\mathbf{p}|}{E+m}\cos(\theta/2) \\
\dfrac{|\mathbf{p}|}{E+m}\sin(\theta/2)e^{i\phi} \\
\cos(\theta/2) \\
\sin(\theta/2)e^{i\phi}
\end{pmatrix}.
\end{align}

Choosing the momentum direction to be the \(z\)-axis (\(\theta = 0\)) givess
\begin{align}
v_\downarrow
=
\sqrt{E+m}
\begin{pmatrix}
\dfrac{|\mathbf{p}|}{E+m} \\
0 \\
1 \\
0
\end{pmatrix}
\equiv
A
\begin{pmatrix}
\tau \xi_R \\
\xi_R
\end{pmatrix},
\end{align}

where
\begin{align}
A = \sqrt{E+m},
\qquad
\tau = \frac{|\mathbf{p}|}{E+m}.
\end{align}

Projecting onto chiral components,
\begin{align}
v_\downarrow = P_L v_\downarrow + P_R v_\downarrow .
\end{align}

The chiral projectors are
\begin{align}
P_R
=
\frac{I_{4\times 4}+\gamma^5}{2},
\qquad
P_L
=
\frac{I_{4\times 4}-\gamma^5}{2}.
\end{align}

Using the definitions of the left- and right-chiral antiparticle states,
\begin{align}
P_L v_R &= v_R, \\
P_R v_L &= v_L ,
\end{align}

one obtains
\begin{align}
v_\downarrow
=
\frac{A}{2}
\left[
(1-\tau)v_R
-
(1+\tau)v_L
\right].
\end{align}

Equivalently,
\begin{align}
v_\downarrow
=
\frac{A}{2}
\left[
\left(1-\frac{|\mathbf{p}|}{E+m}\right)v_R
-
\left(1+\frac{|\mathbf{p}|}{E+m}\right)v_L
\right].
\end{align}

In the limit \(m\to 0\),
\begin{align}
\frac{|\mathbf{p}|}{E+m} \to 1,
\end{align}

so
\begin{align}
v_\downarrow \to -A v_L .
\end{align}

Thus the right-chiral component of a negative-helicity antiparticle vanishes in the massless limit. Since the charged weak interaction couples to the right-chiral antiparticle component, the amplitude is suppressed by the lepton mass.

\paragraph{Pion decay matrix element}

For \(\pi^- \to \ell^- \bar{\nu}_\ell\), the tree-level matrix element is
\begin{align}
\mathcal{M}_{fi}
=
\left[
\frac{g_W}{\sqrt{2}}\frac{1}{2} f_\pi p_\pi^\alpha
\right]
\left[
\frac{g_{\alpha\beta}}{m_W^2}
\right]
\left[
\frac{g_W}{\sqrt{2}}g_\ell
\bar{u}(p_\ell)\gamma^\beta
\frac{1}{2}(1-\gamma^5)v(p_\nu)
\right].
\end{align}

In the pion rest frame, only \(p_\pi^0=m_\pi\) contributes, so
\begin{align}
\mathcal{M}_{fi}
=
\frac{g_W^2 f_\pi g_\ell}{4m_W^2}
m_\pi
\bar{u}(p_\ell)\gamma^0
\frac{1}{2}(1-\gamma^5)v(p_\nu).
\end{align}

Using
\begin{align}
\bar{u}(p_\ell)\gamma^0
=
u^\dagger(p_\ell)\gamma^0\gamma^0
=
u^\dagger(p_\ell),
\end{align}

this becomes
\begin{align}
\mathcal{M}_{fi}
=
\frac{g_W^2 f_\pi g_\ell}{4m_W^2}
m_\pi
u^\dagger(p_\ell)
\frac{1}{2}(1-\gamma^5)v(p_\nu).
\end{align}

For a nearly massless neutrino,
\begin{align}
\frac{1}{2}(1-\gamma^5)v(p_\nu)
=
v_\uparrow(p_\nu).
\end{align}

Therefore,
\begin{align}
\mathcal{M}_{fi}
=
\frac{g_W^2 f_\pi g_\ell}{4m_W^2}
m_\pi
u^\dagger(p_\ell)v_\uparrow(p_\nu).
\end{align}

Choosing the charged lepton to move in the \(+z\)-direction,
\begin{align}
u(p_\ell)
=
u_\uparrow(p_\ell)+u_\downarrow(p_\ell)
=
\sqrt{E_\ell+m_\ell}
\left[
\begin{pmatrix}
1 \\
0 \\
\dfrac{p}{E_\ell+m_\ell} \\
0
\end{pmatrix}
+
\begin{pmatrix}
0 \\
1 \\
0 \\
-\dfrac{p}{E_\ell+m_\ell}
\end{pmatrix}
\right].
\end{align}

The antineutrino spinor is
\begin{align}
v(p_\nu)
=
v_\uparrow(p_\nu)
=
\sqrt{p}
\begin{pmatrix}
1 \\
0 \\
-1 \\
0
\end{pmatrix}.
\end{align}

The negative-helicity charged-lepton contribution vanishes in the spinor product, leaving
\begin{align}
\mathcal{M}_{fi}
=
\frac{g_W^2 f_\pi g_\ell}{4m_W^2}
m_\pi
\sqrt{E_\ell+m_\ell}\sqrt{p}
\left(
1-\frac{p}{E_\ell+m_\ell}
\right).
\end{align}

For \(m_\nu \simeq 0\), the two-body kinematics give
\begin{align}
E_\ell
=
\frac{m_\pi^2+m_\ell^2}{2m_\pi},
\qquad
p
=
\frac{m_\pi^2-m_\ell^2}{2m_\pi}.
\end{align}

Substituting these into the amplitude,
\begin{align}
\mathcal{M}_{fi}
=
\frac{g_W^2 f_\pi g_\ell}{4m_W^2}
m_\pi
\cdot
\frac{m_\pi+m_\ell}{\sqrt{2m_\pi}}
\cdot
\left(
\frac{m_\pi^2-m_\ell^2}{2m_\pi}
\right)^{1/2}
\cdot
\frac{2m_\ell}{m_\pi+m_\ell}.
\end{align}

After simplification,
\begin{align}
\mathcal{M}_{fi}
=
\frac{g_W^2 f_\pi g_\ell}{4m_W^2}
m_\ell
\left(m_\pi^2-m_\ell^2\right)^{1/2}.
\end{align}

Thus,
\begin{align}
|\mathcal{M}_{fi}|^2
=
\left(
\frac{g_W}{2m_W}
\right)^4
f_\pi^2 g_\ell^2
m_\ell^2
\left(m_\pi^2-m_\ell^2\right).
\end{align}

\paragraph{Decay rate}

For a two-body decay,
\begin{align}
\Gamma(\pi \to \ell \bar{\nu}_\ell)
=
\frac{p|\mathcal{M}_{fi}|^2}{8\pi m_\pi^2}.
\end{align}

Using
\begin{align}
p
=
\frac{m_\pi^2-m_\ell^2}{2m_\pi},
\end{align}

the decay rate is
\begin{align}
\Gamma(\pi \to \ell \bar{\nu}_\ell)
=
\frac{f_\pi^2}{16\pi m_\pi^3}
\left(
\frac{g_W}{2m_W}
\right)^4
\left[
m_\ell g_\ell
\left(m_\pi^2-m_\ell^2\right)
\right]^2 .
\end{align}

\paragraph{Tree-level ratio}

Define
\begin{align}
R_{e/\mu}
\equiv
\frac{
\Gamma(\pi^- \to e^- \bar{\nu}_e)
}{
\Gamma(\pi^- \to \mu^- \bar{\nu}_\mu)
}.
\end{align}

Then
\begin{align}
R_{e/\mu}
=
\left(
\frac{g_e}{g_\mu}
\right)^2
\left[
\frac{
m_e(m_\pi^2-m_e^2)
}{
m_\mu(m_\pi^2-m_\mu^2)
}
\right]^2 .
\end{align}

Assuming lepton flavor universality,
\begin{align}
g_e = g_\mu ,
\end{align}

so the tree-level prediction is
\begin{align}
R_{e/\mu}^{\mathrm{tree}}
=
\frac{m_e^2}{m_\mu^2}
\left(
\frac{m_\pi^2-m_e^2}
     {m_\pi^2-m_\mu^2}
\right)^2 .
\end{align}
